Bubble size controls how much mineral a flotation cell recovers, and it is
summarised by the Sauter mean diameter (D32). In practice D32 is still
measured from video largely by hand, which is slow and subjective. This
project built BubbleMetrix, a desktop tool that measures bubble size and
D32 automatically from flotation video, for both the pulp and the froth
zone.

For the pulp, a Mask R-CNN instance segmentation model was trained on
hand-labelled probe footage. Tested on operating conditions it had never
seen, it measured D32 within 1.5\% of careful hand measurement on average
(ratio 0.99 ± 0.03 over 8 conditions), as consistent as a person repeating
their own measurement. For the froth, Cellpose, a general-purpose
segmentation model, was used without any froth-specific training. It found
86\% of hand-outlined bubbles and gave a steady D32 of 2.16 mm, in under a
minute per frame.

BubbleMetrix needs no machine-learning knowledge to use: it measures a
whole video in one step, reports D32 with a 95\% confidence interval, stops
once D32 is known to ±5\%, and can be fine-tuned on new footage when a
quick check shows it is needed.
